\section{Background and Related Work}
\label{sec:background}

This section briefly reviews the foundational concepts and related approaches in static analysis, graph ranking, and automated code retrieval.

\subsection{Static Code Analysis and Graph Representations}

Retrieving code via structure requires extracting relationships from source text without executing it. Standard approaches parse code into an Abstract Syntax Tree (AST), which captures syntactic correctness but can be overwhelmingly detailed for retrieval. In CodeGraph, we utilize Tree-sitter \cite{treesitter} to extract coarse-grained entities (files, classes, functions, methods) rather than full AST nodes.

The structural relationships between these entities (e.g., function calls, module imports, inheritance) are best modeled as a directed graph. Dependency graphs abstract away syntax and expose how information and control flow through a codebase. Graph-based representation has been widely used in software engineering for tasks like vulnerability detection and code comprehension, but translating natural-language queries to these structures remains a core challenge.

\subsection{Graph-Based Ranking via PageRank}

PageRank determines the centrality of a node based on the quality and quantity of its incoming edges. While standard PageRank distributes relevance uniformly across the entire graph, Personalized PageRank (PPR) introduces a bias: the random walk restarts only from a specific set of ``seed'' nodes \cite{page1999pagerank}. This creates a localized probability distribution that ranks nodes based on their proximity and structural connection to the seeds.

For code retrieval, PPR offers a mechanism to combat the semantic gap. If a task description yields a few accurate seed entities, PPR propagates that relevance downstream (callees) and upstream (callers) to rank the neighbourhood most likely to contain the required modification.

\subsection{Related Work in Code Retrieval}

Retrieval for Large Language Models (LLMs) in software engineering typically falls into three broad categories: sparse retrieval, dense retrieval, and graph-based retrieval.

\textbf{Sparse Retrieval:} Sparse retrieval methods, such as BM25, compute relevance based on term overlap between the query and the indexed code tokens. While fast, scalable, and interpretable, sparse methods suffer from extreme lexical sensitivity. A natural language query mentioning ``database compilation'' will not natively match a class named \texttt{DBCompiler} unless specific tokenization rules apply. Thus, pure lexical search struggles when the issue description uses vocabulary that abstracts away from the precise implementation identifiers.

\textbf{Dense Retrieval:} Dense retrieval methods encode queries and code entities into a shared continuous vector space using pre-trained models like CodeBERT or UniXcoder. These methods capture semantic similarity and are more tolerant of paraphrasing. However, dense embeddings compute similarity purely from token sequences---they do not natively understand if a file depends on another or inherits from a base class. Furthermore, dense models are computationally expensive and offer opaque scores without interpretable structural traces.

\textbf{Graph-Based Retrieval:} Graph-based systems explicitly model dependencies (calls, imports, inheritance) and propagate relevance along these paths. Recent systems like RepoGraph, LocAgent, CodexGraph, and LARGER have shown that traversing the graph structure can recover code missed by simple text search. However, existing graph methods typically rely on unranked $k$-hop ego-graph expansion (which dilutes relevance as the radius grows) or LLM-mediated graph queries (which are slow and non-deterministic). The research gap remains: there is a need for a structural retrieval method that is globally ranked, deterministic, and scalable enough to serve as a low-cost agent tool under a strict context window budget.
